\section{Series, products, and exact source checking}
\label{sec:series-products}

\subsection{The binomial-series representation}

For $\Rea x,\Rea y>0$, the source pages give
\begin{equation}
 \B(x,y)=\sum_{n=0}^{\infty}
 \frac{(1-x)_n}{(y+n)n!}.
 \label{eq:beta-series}
\end{equation}
Here $(w)_n$ is the rising factorial.

\begin{theorem}[Effective series formula]\label{thm:beta-series}
Equation \eqref{eq:beta-series} holds as an equality of regular names.  On every certified compact domain
\[
 \Rea x\ge\alpha,\quad \Rea y\ge\beta,\quad |x|,|y|\le R,
\]
there is a primitive-recursive cutoff modulus.
\end{theorem}

\begin{proof}
By symmetry write
\[
 \B(x,y)=\int_0^1 t^{y-1}(1-t)^{x-1}\,dt.
\]
Fix the requested precision.  Remove $t\in[1-\delta,1]$ using the endpoint bound $2\delta^\alpha/\alpha$.  On $[0,1-\delta]$ expand
\[
 (1-t)^{x-1}
 =\sum_{n<N}\frac{(1-x)_n}{n!}t^n+R_N(t).
\]
This is the finite Taylor formula for $\exp((x-1)\log(1-t))$.  The elementary exponential remainder and the bound $|\log(1-t)|\le\log(1/\delta)$ give an explicit factorial cutoff $N$.  Integrating the finite polynomial gives
\[
 \sum_{n<N}\frac{(1-x)_n}{n!}\int_0^{1-\delta}t^{y+n-1}\,dt.
\]
Replacing the upper endpoint $1-\delta$ by $1$ contributes another explicitly bounded endpoint term, and
$\int_0^1t^{y+n-1}\,dt=(y+n)^{-1}$.  The finite algebra is exact after guarded inversion because $\Rea(y+n)\ge\beta+n$.  All cutoffs are primitive recursive in the domain data and precision.
\end{proof}

For the incomplete beta function the same proof gives the more useful uniformly geometric series on $0\le z\le1-\delta$:
\[
 \B_z(a,b)=
 \sum_{n=0}^{\infty}
 \frac{(1-b)_n}{(a+n)n!}z^{a+n}.
\]
This will be used in \cref{sec:incomplete}.

\subsection{A displayed symmetric series is false as written}

The Wikipedia page currently displays
\[
 \B(z,z)=2\sum_{n\ge0}
 \frac{(2z+n-1)_n(-1)^n}{(z+n)n!}
 =\lim_{r\to1^-}2\sum_{n\ge0}
 \frac{(-2z)_n r^n}{(z+n)n!},
 \tag{S}
\]
while explicitly declaring $(\cdot)_n$ to be the rising factorial.

\begin{warning}[Counterexample at $z=1$]\label{warn:false-series}
Formula (S) is false with the stated rising-factorial convention.
\end{warning}

\begin{proof}
For $z=1$, the magnitude of the $n$th term in the first series is
\[
 2\frac{(n+1)_n}{(n+1)n!}
 =2\frac{(2n)!}{(n+1)(n!)^2},
\]
which does not tend to zero; the series diverges.  In the second series, $(-2)_n$ vanishes for $n\ge3$, so the Abel limit is
\[
 2\left(1-1+\frac{2}{3\cdot2}\right)
 =\frac23,
\]
not $\B(1,1)=1$.  Thus neither equality can be used in a proof packet without correcting the source formula.
\end{proof}

The exact-rational checker in the archive records this counterexample.  We do not guess the intended correction.

\subsection{The Euler product for Beta}

The source pages give two algebraically equivalent products:
\begin{align}
 \B(x,y)
 &=\frac{x+y}{xy}
 \prod_{k=1}^{\infty}
 \frac{1+(x+y)/k}{(1+x/k)(1+y/k)},
 \label{eq:beta-product-1}\\
 &=\frac{x+y}{xy}
 \prod_{k=1}^{\infty}
 \left(1+\frac{xy}{k(k+x+y)}\right)^{-1}.
 \label{eq:beta-product-2}
\end{align}
The equivalence of their $k$th factors is the exact identity
\[
 \frac{(1+x/k)(1+y/k)}{1+(x+y)/k}
 =1+\frac{xy}{k(k+x+y)}.
\]

\begin{theorem}[Effective infinite product]\label{thm:beta-product}
On a certified compact domain separated from the poles and from $x=0$, $y=0$, and $x+y=0$, the finite products in
\eqref{eq:beta-product-1}--\eqref{eq:beta-product-2} converge to $\B(x,y)$ with a primitive-recursive modulus.
For positive real $x,y$, the tail after $N$ satisfies
\[
 1\le
 \prod_{k>N}\left(1+\frac{xy}{k(k+x+y)}\right)
 \le\exp(xy/N).
\]
\end{theorem}

\begin{proof}
Insert the effective Weierstrass products for $\Gamma(x)$, $\Gamma(y)$, and $\Gamma(x+y)$ into the beta--gamma quotient.  The exponential normalization factors cancel exactly, leaving \eqref{eq:beta-product-1}.  This cancellation occurs already at each common finite cutoff, so only the Gamma product remainder needs to be bounded.

For positive $x,y$, put $u_k=xy/[k(k+x+y)]$.  Then $u_k\ge0$ and
\[
 \sum_{k>N}u_k\le xy\sum_{k>N}\frac1{k^2}\le\frac{xy}{N}.
\]
The finite inequality $\prod(1+u_k)\le\exp(\sum u_k)$ follows by multiplying the finite inequalities $1+u\le e^u$.  The lower bound is immediate.  For complex parameters use $|u_k|\le C/k^2$ once $k$ exceeds the fixed radius, and the finite telescoping product estimate gives a tail bounded by a primitive-recursive multiple of $1/N$.
\end{proof}

For positive $x,y$ this product also gives monotone rational upper and lower approximations to Beta after the prefactor is bounded.  It is therefore an alternative evaluator to the integral name.

\subsection{Finite products before limits}

A useful formalization policy is to prove the identity
\[
 \frac{x+y}{xy}
 \prod_{k=1}^{N}
 \frac{1+(x+y)/k}{(1+x/k)(1+y/k)}
 =
 \frac{x+y}{xy}\,
 \frac{N!\,(x+y+1)_N}{(x+1)_N(y+1)_N}
\]
only after clearing all denominators.  A simpler equivalent form is obtained directly from the Euler Gamma approximants.  The exact choice is not mathematically important; what matters is that the product regrouping is a quantifier-free induction invariant and the analytic content is isolated in an explicit $O(1/N)$ remainder.

\subsection{Products and logarithms}

On positive real domains,
\[
 \log\B(x,y)=\log(x+y)-\log x-\log y
 -\sum_{k\ge1}\log\left(1+\frac{xy}{k(k+x+y)}\right).
\]
Since $0\le\log(1+u)\le u$, the same $xy/N$ bound controls the logarithmic tail.  Finite differentiation of this logarithmic product supplies an independent route to the digamma formulas in \cref{sec:derivatives}; the Gamma quotient is shorter, but the product route may be preferable in an object-theory development because every denominator is visibly separated on a positive compact domain.
